\documentclass[nooutcomes]{ximera}

\usepackage{gensymb}
\usepackage{tabularx}
\usepackage{mdframed}
\usepackage{pdfpages}
%\usepackage{chngcntr}

\let\problem\relax
\let\endproblem\relax

\newcommand{\property}[2]{#1#2}




\newtheoremstyle{SlantTheorem}{\topsep}{\fill}%%% space between body and thm
 {\slshape}                      %%% Thm body font
 {}                              %%% Indent amount (empty = no indent)
 {\bfseries\sffamily}            %%% Thm head font
 {}                              %%% Punctuation after thm head
 {3ex}                           %%% Space after thm head
 {\thmname{#1}\thmnumber{ #2}\thmnote{ \bfseries(#3)}} %%% Thm head spec
\theoremstyle{SlantTheorem}
\newtheorem{problem}{Problem}[]

%\counterwithin*{problem}{section}



%%%%%%%%%%%%%%%%%%%%%%%%%%%%Jenny's code%%%%%%%%%%%%%%%%%%%%

%%% Solution environment
%\newenvironment{solution}{
%\ifhandout\setbox0\vbox\bgroup\else
%\begin{trivlist}\item[\hskip \labelsep\small\itshape\bfseries Solution\hspace{2ex}]
%\par\noindent\upshape\small
%\fi}
%{\ifhandout\egroup\else
%\end{trivlist}
%\fi}
%
%
%%% instructorIntro environment
%\ifhandout
%\newenvironment{instructorIntro}[1][false]%
%{%
%\def\givenatend{\boolean{#1}}\ifthenelse{\boolean{#1}}{\begin{trivlist}\item}{\setbox0\vbox\bgroup}{}
%}
%{%
%\ifthenelse{\givenatend}{\end{trivlist}}{\egroup}{}
%}
%\else
%\newenvironment{instructorIntro}[1][false]%
%{%
%  \ifthenelse{\boolean{#1}}{\begin{trivlist}\item[\hskip \labelsep\bfseries Instructor Notes:\hspace{2ex}]}
%{\begin{trivlist}\item[\hskip \labelsep\bfseries Instructor Notes:\hspace{2ex}]}
%{}
%}
%% %% line at the bottom} 
%{\end{trivlist}\par\addvspace{.5ex}\nobreak\noindent\hung} 
%\fi
%
%


\let\instructorNotes\relax
\let\endinstructorNotes\relax
%%% instructorNotes environment
\ifhandout
\newenvironment{instructorNotes}[1][false]%
{%
\def\givenatend{\boolean{#1}}\ifthenelse{\boolean{#1}}{\begin{trivlist}\item}{\setbox0\vbox\bgroup}{}
}
{%
\ifthenelse{\givenatend}{\end{trivlist}}{\egroup}{}
}
\else
\newenvironment{instructorNotes}[1][false]%
{%
  \ifthenelse{\boolean{#1}}{\begin{trivlist}\item[\hskip \labelsep\bfseries {\Large Instructor Notes: \\} \hspace{\textwidth} ]}
{\begin{trivlist}\item[\hskip \labelsep\bfseries {\Large Instructor Notes: \\} \hspace{\textwidth} ]}
{}
}
{\end{trivlist}}
\fi


%% Suggested Timing
\newcommand{\timing}[1]{{\bf Suggested Timing: \hspace{2ex}} #1}




\hypersetup{
    colorlinks=true,       % false: boxed links; true: colored links
    linkcolor=blue,          % color of internal links (change box color with linkbordercolor)
    citecolor=green,        % color of links to bibliography
    filecolor=magenta,      % color of file links
    urlcolor=cyan           % color of external links
}

\title{Instructor Overview: Math for Elementary Teachers Activities}
\author{Vic Ferdinand, Betsy McNeal, Jenny Sheldon}

\begin{document}

\begin{abstract}

\end{abstract}\maketitle

\includegraphics[height=2in]{OOECLogo.png}

In these pages, we have gathered a collection of the activities used at The Ohio State University in a 
two-semester sequence (Math 1125 and 1126) of mathematics content courses for preservice elementary 
school teachers.  This sequence is also taken as the mathematics requirement for students preparing 
to teach Special Education in grades K-12 and for those preparing to teach subjects other than 
mathematics in middle school.

This collection of activities is not intended to be a full course or a textbook, or to cover all of the 
content required for preservice elementary school teachers.  Rather, the activities are intended to 
stand either alone or in small units, so that they can be reorganized or used in part, as an 
instructor deems most appropriate.  Within each file download, we have provided first the activity 
itself, followed by Instructor's Notes.  These Instructor's Notes are designed to give context about 
how we use the activity at Ohio State, as well as what content is assumed before starting the 
activity. At Ohio State, the courses are each five credit 
hours, and run five days per week.  Many of the activities are designed to take the whole class 
period, or 55 minutes, while others use approximately half the class.  
The Instructor's Notes often give a suggested amount of time for the 
activity, but reference our 55-minute ``class period''.

The source code for the activities is available on GitHub at 
\url{https://github.com/mooculus/mathForTeachers}.  Familiarity with \LaTeX as well as Ohio State's
Ximera document class (\url{https://ximera.osu.edu/}) is recommended for editing the source code.
One open source textbook we discovered while working on this project is ``Mathematics for Elementary 
Teachers'' by Michelle Manes, available at 
\link[Open Textbook Library]{https://open.umn.edu/opentextbooks/textbooks/mathematics-for-elementary-teachers}.




At Ohio State, we have several main goals for our course.  First, we aim to have students refresh and deepen their 
understanding of elementary school mathematical procedures and concepts.  Although most of our 
students have low interest in mathematics and lack confidence in this subject area, they often 
THINK they already know elementary school math and tend to respond to a request 
for further explanation with either ``that's what I learned in school'' or ``I don't know''.  Throughout the 
course, we try to challenge their starting assumptions and provoke deeper reasoning by selecting 
text problems and supplements that cause our students to approach these basic ideas from fresh 
points of view.  We constantly ask students why a procedure works as it does and what it 
means.  We do not accept a formula as an explanation.  We require explanations that are
based on the meaning of a number (e.g., a fraction) and the meaning of an operation (e.g., 
multiplication as a number of groups with the same number of objects in each group).

One theme we see in our activities is an intention to ``make the familiar strange''.  To us, this means 
we would like to take mathematical concepts that are familiar to students, but approach them from 
angles which make the problems or concepts unfamiliar.  When students recognize a problem, we have 
found that they are likely to try to apply a ``rule'' which they remember from school, often 
without much thought as to whether that rule applies or is appropriate.  When we take away the 
familiar contexts, students seem more willing to use definitions and think through concepts instead
of just applying a rule.

Our second goal is to encourage a positive attitude towards mathematics and an enriched view 
of what mathematics is all about.  We hope to move students towards an understanding of mathematics
as a process of reasoning and away from a view of math as a body of 
rules and facts to be memorized. Such a view of mathematics is important for teachers to make a productive and 
creative classroom. A teacher who understands the intent of mathematics as a way of 
thinking will encourage her students to engage more deeply in that way of thinking.  

Both of these goals are addressed in our emphasis on explanation.  The more 
advanced students are challenged to view the focus on explanations and grade appropriate justifications as building toward the more rigorous proofs they experienced in high school.  
Students' explanations must follow directly and logically from the context in which we are working.  
All class meetings (lecture and recitation) at OSU are designed to get students thinking about, talking 
about, writing, and representing mathematical ideas.  This encourages deeper and more connected 
understanding of the mathematics on their part as well as relating the learning to their preparation as future teachers.  

Our class meetings are frequently composed of deep discussion about a handful of problems.  Since we also require 
a textbook (currently ``Mathematics for Elementary Teachers, 5ed'' by Sybilla Beckmann), we rely on 
students reading the text after class to help with background information or other details.  The 
classroom activities were designed with such background reading in mind.

As a sort of ``capstone'' experience, we designed a unit of projects for the end of the second 
course.  These projects help us to expose students to some more sophisticated content without 
having to require mastery of this content by all students.  They also help build problem-solving skills and 
encourage independent learning.  In the ``Instructor's Notes'' with some activities in the second 
course, you'll see references to places some of the content can be left for the projects at the end. 
These are indications that the content is more advanced and could equally well be skipped.

Some instructors might read our activities and initially think they are easy or superficial.  The purpose of these notes is to illustrate how these activities can be used to probe the depths of each topic in the traditional content for elementary mathematics teachers.  Our descriptions are drawn from our actual experiences of teaching many future teachers over many years.

Finally, we have more detailed solutions and Instructor's Notes for many of these activities than those 
which are publicly posted here.  We are happy to share those privately as well as answer any 
questions you may have about the activities, calendar, or other aspects of the course.  Please feel 
free to contact Vic Ferdinand, Betsy McNeal, or Jenny Sheldon at Ohio State for more information.

We would like to give our deepest thanks to Sybilla Beckmann and acknowledge the significant role her textbook, \emph{Mathematics for Elementary Teachers with Activities}, has played in the success of our program.  We hope that we have properly separated our own problems and activities from those in her textbook. We have followed her through five editions and continue to find this a supportive and readable text for our students and, more important, one that is consistent with our pedagogical goals. We are able to encourage messy discussions based on our students' thinking because we know that our students are going to read her clear text after our discussions.  

Finally, we would like to thank the Ohio Department of Higher Education, which in 2017 awarded a 1.3 million dollar grant to North Central State College and two partner institutions — The Ohio State University and Ohio Dominican University — with the goal of creating new and curating existing open materials for over twenty courses.  These materials were put together in this format as a result of this grant.



\end{document}
